%\documentclass{Math_Note}
%\setlength{\parindent}{2em}

%\title{\textbf{Mathmatical Analysis Zorich Comprehensions Chapter 5}}
%\author{Buce-Ithon}
%\newdateformat{mydate}{\twodigit{\THEDAY}{ }\shortmonthname[\THEMONTH], \THEYEAR}
%\date{\today}

%\begin{document}
% Title page
%\maketitle

% Content
%\newpage
%\tableofcontents
%\newpage

% \setcounter{section}{-1}
\newpage
\section{Chapter 5.Differential Calculus}
\subsection{Comprehensions}
1st part: Introduce the definition of functions differentiable at a point and related concepts: differential, derivative and increment. 
Besides, it also includes the geometrical meaning of them and some calculation examples. 
After basic definitions, we are given the basic rules of differentiation, including: arithmetic operations, differentiation of composite 
functions, differentiation of inverse functions, and higer order derivatives (Lebniz's formula). And also derivatives of basic elementary 
functions. \\
2nd part: Then we dive into some basic theorems of differential calculus: /All of beging-Most basic/ \underline{Fermat's Lemma}, /Dev on 
basic-3 Mean Value Theorems/ \underline{Rolle's Theorem} $\to$ \underline{Lagrange's Mean Value Theorem} $\to$ \underline{Cauchy's Mean 
Value} \underline{Theorem} (due to last 2 of them associated the increment and derivative, they're also called Theorem of Lagrange and 
Cauchy on finite increments). Finally, we're given the most important theorem in differential calculus: \underline{Taylor's Formula} (the 
reminder term could be obtained using Cauchy's Mean Value Theorem or Theorem of Cauchy on finite increments). \\
3nd part: It's the last part of this chapter, which focuses on the applications (analyze properties of functions) of differential calculus, 
including: monotoncity, interior extremum, convacity, and graph sketching. And the well-known L'H\^opital's Rule. Remains lots of examples 
in Problems on Natural Science. \\
DLC: The special (which could be called overall) point is that Zorich left a specific section for complex numbers and functions, even including 
analyticity (or holomorphism). 
\newline\newline
\fbox{\begin{minipage}{0.95\textwidth}
\textbf{Definition: } /differential/ $\underbrace{\underbrace{f(x+h)-f(x)}=\overbrace{A(x)\underbrace{h}}^{\text{differential}}}_{\text{increments}}+
\underbrace{o(a;h)}_{o(h),\ h\to 0}$, $x \in E$ 
/derivative/ $f'(x) = A(x) = \lim_{h\to 0}\frac{f(x+h)-f(x)}{h}$
\end{minipage}}
\newline\newline
All of beging-Most basic \\
\fbox{\begin{minipage}{0.95\textwidth}
\begin{thm}
/Fermat's Lemma/ \\
$f : E \to \mathbb{R}$ is differentiable at an interior extremum $x_0 \in E$ $\implies$  $f'(x_0) = 0$.
\end{thm}
\end{minipage}}
\newline\newline
3 Mean Value Theorems \\
\fbox{\begin{minipage}{0.95\textwidth}
\begin{thm}
/Rolle's theorem/ \\
$f : [a, b] \to \mathbb{R}$ is continuous on $[a, b]$ and differentiable on $(a, b)$ and $f(a) = f(b)$, then $\exists \xi \in (a, b)$ s.t. $f'(\xi) = 0$.
\end{thm}
\end{minipage}}
\newline\newline
\fbox{\begin{minipage}{0.95\textwidth}
\begin{thm}
/Lagrange's Mean Value Theorem/ \\
$f : [a, b] \to \mathbb{R}$ is continuous on $[a, b]$ and differentiable on $(a, b)$, then $\exists \xi \in (a, b)$ s.t. $f(b) - f(a) = f'(\xi)(b - a)$.
\end{thm}
\end{minipage}}
\newline
\fbox{\begin{minipage}{0.95\textwidth}
\begin{thm}
/Cauchy's Mean Value Theorem/ \\
Let $x = x(t)$ and $y = y(t)$ be continuous on $[\alpha, \beta]$ and differentiable on $(\alpha, \beta)$.
Then $\exists \tau \in [\alpha, \beta]$ s.t. $x'(\tau)(y(\beta) - y(\alpha)) = y'(\tau)(x(\beta) - x(\alpha))$. \\
If in addition $x'(t) \neq 0$, $\forall t \in [\alpha, \beta]$, then $x(\alpha) \neq x(\beta)$ and $\frac{y(\beta) - y(\alpha)}{x(\beta) - x(\alpha)} 
= \frac{y'(\tau)}{x'(\tau)}$.
\end{thm}
\end{minipage}}
\newline\newline
Master in differential calculus \\
\fbox{\begin{minipage}{0.95\textwidth}
\begin{thm}
/Taylor's Formula/ \\
$$F(a) = f(x_0) + f'(x_0)(x - x_0) + \cdots + \frac{f^{(n)}(x_0)}{n!}(x - x_0)^n + r_n(x_0, x)$$
\end{thm}
\end{minipage}}
\newline\newline
\fbox{\begin{minipage}{0.95\textwidth}
\textbf{Remark: } /Reminder term/ \\
$$r_n(x_0;x) = \frac{\varphi(x) - \varphi(x_0)}{\varphi'(\xi)n!} f^{(n+1)}(\xi)(x-\xi)^n.$$
or $$\frac{r_n(x_0;x)}{\varphi(x) - \varphi(x_0)}=\frac{f^{(n+1)}(\xi)(x-\xi)^n}{\varphi'(\xi)n!} \ \ \text{/Cauchy's MVT/}$$
/Cauchy’s formula of the remainder term/ \\
$$r_n(x_0; x) = \frac{1}{n!} f^{(n+1)}(\xi)(x - \xi)^n(x - x_0).$$
A particularly elegant formula results if we set $\varphi(t) = (x - t)^{n+1}$ \\
/The Lagrange form of the remainder/ \\
$$r_n(x_0; x) = \frac{1}{(n + 1)!} f^{(n+1)}(\xi)(x - x_0)^{n+1}.$$
We remark that when $x_0 = 0$ Taylor's formula is often called MacLaurin's formula.
\end{minipage}}

\newpage
\subsection{Problems Sets}
\begin{prb}
    (C5S1Q5) /Continuous, but nowhere derivative/ \\
    Set $\Psi_0(x) = \begin{cases} x, & \text{if } 0 \leq x \leq \frac{1}{2}, \\ 1 - x, & \text{if } \frac{1}{2} \leq x \leq 1, \end{cases}$ and extend this 
    function to the entire real line so as to have period 1. We denote the extended function by $\varphi_0$. Further, let $\varphi_n(x) = \frac{1}{4^n} 
    \varphi_0(4^n x)$. The function $\varphi_n$ has period $4^{-n}$ and a derivative equal to +1 or $-1$ everywhere except at the points $x = \frac{k}{2^{n+1}}, 
    n \in \mathbb{Z}$. Let $f(x) = \sum_{n=1}^\infty \varphi_n(x)$. \\
    Show that the function $f$ is defined and continuous on $\mathbb{R}$, but does not have a derivative at any point.
    (This example is due to the well-known Dutch mathematician B. L. van der Waerden (1903–1996). The first examples of continuous functions having no derivatives 
    were constructed by Bolzano (1830) and Weierstrass (1860).)
\end{prb}
\fbox{\begin{minipage}{0.95\textwidth}
\begin{sol*}
1. Using Weierstrass M-test (uniform convergence of series of functions) and the uniform limit theorem (the uniform limit of any sequence of continuous functions 
is continuous), we can show that $f \in C(\mathbb{R})$. \\
2. Just consider $\lim_{n \to \infty} \frac{f(x\pm 4^{-n})-f(x)}{4^{-n}}=\lim_{n \to \infty} \frac{\sum_{k=0}^{\infty}\varphi_k(x\pm 4^{-n})-\sum_{k}\varphi_k(x)}
{4^{-n}}$, and using the periodicity of $\varphi_k$.  
\end{sol*}
\end{minipage}}
\newline\newline
\fbox{\begin{minipage}{0.95\textwidth}
\textcolor{cyan}{/Dev/} \textbf{More step problem:} \\
Prove that the nowhere differential functions are dense in $C(\mathbb{R})$.
\end{minipage}}

\begin{prb}
    (C5S2Q1) \\
    Let $\alpha_0, \alpha_1, \dots , \alpha_n$ be given real numbers. Find a polynomial $P_n(x)$ of degree $n$ having the derivatives $P_n^{(k)}(x_0) = \alpha_k, 
    k = 0, 1, \dots , n$, at a given point $x_0 \in \mathbb{R}$.
\end{prb}
\fbox{\begin{minipage}{0.95\textwidth}
\begin{qsol*}
$P_n(x)=\sum_{k=0}^{n} \alpha_k(x-x_0)^k$ 
\end{qsol*}
\end{minipage}}

\begin{prb}
    (C5S2Q2) 
\end{prb}
\fbox{\begin{minipage}{0.95\textwidth}
\begin{qsol*}
\begin{enumerate}
    \item[a)] $f(x)=\begin{cases} e^{-\frac{1}{x^2}} & x\neq 0 \\ 0 & x=0 \end{cases}$,$f'(x)=\begin{cases} \frac{2e^{-\frac{1}{x^2}}}{x^3} 
        & x\neq 0 \\ 0 & x=0 \end{cases}$, $f^{(n)}(x)=\lim_{x\to\infty}\frac{f^{(n-1)}(x)-0}{x-0}=[\text{Induction}]=\lim_{x\to\infty}\frac{R_n(x)}
        {e^{\frac{1}{x^2}}}=0$. 
    \item[b)] $f(x)=\begin{cases} x^2\sin(\frac{1}{x}) & x\neq0 \\ 0 & x=0 \end{cases}$, $f(x)=\begin{cases} 2x\sin(\frac{1}{x})-\cos(\frac{1}{x}) & x\neq0 
        \\ 0 & x=0 \end{cases}$ 
    \item[c)] $f(x)=\begin{cases} e^{-\frac{1}{(1+x)^2}}\cdot e^{-\frac{1}{(1-x)^2}} & -1<x<1 \\ 0 & 1\leq|x| \end{cases}$ 
\end{enumerate}
\end{qsol*}
\end{minipage}}

\begin{prb}
    (C5S2Q3) \\
    $$\left[ {x^{n - 1} f\left( {\frac{1}{x}} \right)} \right]^{(n)} = \frac{{( - 1)^n }}{{x^{n + 1} }}f^{(n)} \left( {\frac{1}{x}} \right)$$
\end{prb}
\fbox{\begin{minipage}{0.95\textwidth}
\begin{qsol*}
$[$Induction$]$. Assume that the base case $n=1$ is proved. \\
Assuming that $\frac{d^n}{dx^n}\left[x^{n-1}f\left(\frac{1}{x}\right)\right]=\frac{(-1)^n}{x^{n+1}}f^{(n)}\left(\frac{1}{x}\right)$, we have $$\frac{d^{n+1}}
{dx^{n+1}}\left[x^{n}f\left(\frac{1}{x}\right)\right]=\frac{d^n}{dx^n}\left[nx^{n-1}f\left(\frac{1}{x}\right)-x^{n-2}f'\left(\frac{1}{x}\right)\right]$$ $$
\underbrace{=\frac{n(-1)^{n}}{n^{n+1}}f^{(n)}\left(\frac{1}{x}\right)}_\text{by asumption}-\frac{d}{dx}\left(\frac{d^{n-1}}{dx^{n-1}}\left[x^{n-2}f'\left(
\frac{1}{x}\right)\right]\right)$$ $$=\frac{n(-1)^{n}}{n^{n+1}}f^{(n)}\left(\frac{1}{x}\right)-\frac{d}{dx}\underbrace{\left[\frac{(-1)^{n-1}}{x^n}f^{(n)}\left(
\frac{1}{x}\right)\right]}_\text{by assumption}$$ $$=\frac{n(-1)^{n}}{n^{n+1}}f^{(n)}\left(\frac{1}{x}\right)+\frac{n(-1)^{n-1}}{x^{n+1}}f^{(n)}\left(\frac{1}{x}\right)
+\frac{(-1)^{n+1}}{x^{n+2}}f^{(n+1)}\left(\frac{1}{x}\right)$$ $$=\frac{(-1)^{n+1}}{x^{n+2}}f^{(n+1)}\left(\frac{1}{x}\right)\square$$
And that about wraps it up. \textit{"from AOPS"}
\end{qsol*}
\end{minipage}}

\begin{prb}
    (C5S2Q4) \\
    Let $f$ be a differentiable function on $\mathbb{R}$. Show that
    \begin{enumerate}
        \item[a)] if $f$ is an even function, then $f'$ is an odd function;
        \item[b)] if $f$ is an odd function, then $f'$ is an even function;
        \item[b)] ($f'$ is odd) $\Leftrightarrow$ ($f$ is even).
    \end{enumerate}
\end{prb}
\fbox{\begin{minipage}{0.95\textwidth}
\begin{qsol*}
By definition.
\end{qsol*}
\end{minipage}}

\begin{prb}
    (C5S2Q5) \\
    Show that
    \begin{enumerate}
        \item[a)] the function $f(x)$ is differentiable at the point $x_0$ if and only if $f(x) - f(x_0) = \varphi(x)(x - x_0)$, where $\varphi(x)$ is a function that 
        is continuous at $x_0$ (and in that case $\varphi(x_0) = f'(x_0)$);
        \item[b)] if $f(x) - f(x_0) = \varphi(x)(x - x_0)$ and $\varphi \in C^{(n-1)} \left( U(x_0) \right)$, where $U(x_0)$ is a neighborhood of $x_0$, then $f(x)$ has a 
        derivative $(f^{(n)}(x_0))$ of order $n$
    \end{enumerate}
\end{prb}
\fbox{\begin{minipage}{0.95\textwidth}
\begin{qsol*}
By definition.
\end{qsol*}
\end{minipage}}

\begin{prb}
    (C5S2Q6) \\
    Give an example showing that the assumption that $f^{-1}$ be continuous at $y_0$ cannot be omitted from following theorem: \\
    \begin{thm}
        (The derivative of an inverse function). Let the functions $f : X \to Y$ and $f^{-1} : Y \to X$ be mutually inverse and continuous at points $x_0 \in X$ and 
        $f(x_0) = y_0 \in Y$ respectively. If $f$ is differentiable at $x_0$ and $f'(x_0) \neq 0$, then $f^{-1}$ is also differentiable at the point $y_0$, and 
        $(f^{-1})'(y_0) = (f'(x_0))^{-1}$.
    \end{thm}
\end{prb}
\fbox{\begin{minipage}{0.95\textwidth}
\begin{qsol*}
Let $f(x)=\begin{cases} x+x^2\sin\frac{1}{x}, & x\neq 0 \\ 0, & x=0 \end{cases}$, after calculating, we could find $f$ is continuous and differentiable, but $f^{-1}$ 
doesn't even exist. 
\end{qsol*}
\end{minipage}}

\begin{prb}
    (C5S3Q1) \\
    Choose numbers $a$ and $b$ so that the function $f(x) = \cos x - \frac{1+ax^2}{1+bx^2}$ is an infinitesimal of highest possible order as $x \to 0$.
\end{prb}
\fbox{\begin{minipage}{0.95\textwidth}
\begin{qsol*}
$b=0,\ a=-\frac{1}{2}$ ($\cos x = 1-\frac{1}{2!}x^2+\frac{1}{4!}x^4+\mathcal{O}(x^6)$)
\end{qsol*}
\end{minipage}}

\begin{prb}
    (C5S3Q2) 
    $$\lim_{x \to \infty} x\left[\frac{1}{e}-(\frac{x}{x+1})^x\right]$$
\end{prb}
\fbox{\begin{minipage}{0.95\textwidth}
\begin{sol*}
\begin{equation*}
\begin{split}
    \lim_{x \to \infty} x&\left[\frac{1}{e}-(\frac{x}{x+1})^x\right] =\lim_{x \to \infty} x\left[\frac{1}{e}-e^{x\ln(\frac{x}{x+1})}\right] \ (\ln\frac{x}{x+1}=\ln(1-\frac{1}{x+1})) \\
    =&\lim_{x \to \infty} x\left[e^{-1}-e^{x(-\frac{1}{x+1}-\frac{1}{2(x+1)^2}+\mathcal{O}(\frac{1}{x^3}))}\right] \\
    =&\lim_{x \to \infty} x\left[e^{-1}-e^{-\frac{x}{x+1}-\frac{x}{2(x+1)^2}+\mathcal{O}(\frac{1}{x^2})}\right] \\
    =&\lim_{x \to \infty} x\left[e^{-1}-e^{-\frac{x}{x+1}}(1-\frac{x}{2(x+1)^2}+\mathcal{O}(\frac{1}{x^2}))\right] \\
    =&\lim_{x \to \infty} e^{-1} \left[x-xe^{\frac{1}{x+1}}+\frac{x^2}{2(x+1)^2}e^{\frac{1}{x+1}}+\mathcal{O}(\frac{1}{x})\right] \\
    =&\lim_{x \to \infty} e^{-1} \left[x(1-(1+\frac{1}{x+1}+\mathcal{O}(\frac{1}{x^2})))+\frac{x^2}{2(x+1)^2}e^{\frac{1}{x+1}}+\mathcal{O}(\frac{1}{x})\right] \\
    =&\lim_{x \to \infty} e^{-1} \left[\frac{x}{x+1}+\frac{x^2}{2(x+1)^2}(1+o(1))+\mathcal{O}(\frac{1}{x})\right] \\
    =&\frac{3}{2}e^{-1}
\end{split}
\end{equation*}
\end{sol*}
\end{minipage}}

\begin{prb}
    (C5S3Q3) 
\end{prb}
\fbox{\begin{minipage}{0.95\textwidth}
\begin{qsol*}
    $e^x=1+x+\frac{1}{2!}x^2+\frac{1}{3!}x^3+\frac{1}{4!}x^4+\cdots$  
\end{qsol*}
\end{minipage}}

\begin{prb}
    (C5S3Q4) \\
    Let $f$ be a function that is infinitelt differentiable at $0$. Show that: 
    \begin{enumerate}
        \item[a)] if $f$ is even, then its Taylor series at $0$ contains only even powers of $x$;
        \item[b)] if $f$ is odd, then its Taylor series at $0$ contains only odd powers of $x$.
    \end{enumerate}
\end{prb}
\fbox{\begin{minipage}{0.95\textwidth}
\begin{qsol*}
    The parity of $f$ and $f^{\prime}$ is opposite. \& When $f(x)=-f(-x)$, $f(0)=0$. 
\end{qsol*}
\end{minipage}}

\begin{prb}
    (C5S3Q5) \\
    Show that if $f \in C^{(\infty)}[-1, 1]$ and $f^{(n)}(0) = 0$ for $n = 0, 1, 2, \ldots$, and there $\exists$ a number $C$ s.t. $\sup_{x \in [-1, 1]} |f^{(n)}(x)| 
    \leq n! C$ for $n \in \mathbb{N}$, then $f = 0$ on $[-1, 1]$. 
\end{prb}
\fbox{\begin{minipage}{0.95\textwidth}
\begin{qsol*}
    Taylor's formula: on $x=0$, $|f(x)| = |\frac{f^{(n+1)}(\varepsilon_n)}{(n+1)!}x^{n+1}| \leq |\frac{(n+1)!C}{(n+1)!}x^{n+1}| = |Cx^{n+1}| = \mathcal{O}(x^n)$. Since 
    $f(0)=0$, and continuity, $f(x)\equiv 0$. 
\end{qsol*}
\end{minipage}}

\begin{prb}
    (C5S3Q6) \\
    Let $f \in C^{(n)}(- 1, 1)$ and $\sup\limits_{-1 < x < 1} |f(x)| \leq 1$. Let $m_k(I) = \inf\limits_{x \in I} |f^{(k)}(x)|$, where $I$ is an interval contained in 
    $(- 1, 1)$. Show that 
    \begin{enumerate}
        \item[a)] if $I$ is partitioned into three successive intervals $I_1, I_2,$ and $I_3$ and $\mu$ is the length of $I_2$, then $m_k(I) \leq \frac{1}{\mu} \left( 
            m_{k-1}(I_1) + m_{k-1}(I_3) \right);$
        \item[b)] if $I$ has length $\lambda$, then $m_k(I) \leq \frac{2^{k(k+1)/2} k^k}{\lambda^k};$
        \item[c)] there exists a number $\alpha_n$ depending only on $n$ such that if $|f'(0)| \geq \alpha_n$, then the equation $f^{(n)}(x) = 0$ has at least $n - 1$ 
            distinct roots in $(- 1, 1)$.
    \end{enumerate}
    Hint: In part b) use part a) and mathematical induction; in c) use a) and prove by induction that there exists a sequence $x_{k_1} < x_{k_2} < \cdots < x_{k_k}$ of 
    points of the open interval $(- 1, 1)$ s.t. $f^{(k)}(x_{k_i}) \cdot f^{(k)}(x_{k_{i+1}}) < 0$ for $1 \leq i \leq k - 1$.
\end{prb}
\fbox{\begin{minipage}{0.95\textwidth}
\begin{sol*}
\begin{enumerate}
    \item[a)] (Lagrange MVT) \\
    $m_k(I)=\inf\limits_{x \in I \subset (-1,1)}|f^{(k)}(x)| \leq m_k(I_2)=\inf\limits_{x \in I_2 \subset I}|f^{(k)}(x)| \leq \forall |f^k(x)| \in \{|f^k(x)|:x \in I_2\} 
    =[|f^k(\xi)|,\xi \in I_2]=\frac{\inf\limits_{x \in I_1 \subset I}|f^{(k-1)}(x)| - \inf\limits_{x \in I_3 \subset I}|f^{(k-1)}(x)|}{x_1-x_2} \leq \frac{m_{k-1}(I_1)+m_{k-1}(I_3)}{\mu}$ 
    \item[b)] $m_k(I)=\inf\limits_{x \in I}|f^{(k)}(x)|,\ \mu(I)=\lambda \leq 2$ \\
        $m_k(I)\leq \frac{1}{\mu}(m_{k-1}(I_1)+m_{k-1}(I_3)) = [\mu=\mu(I_2)\triangleq\frac{\lambda}{2}, \mu(I_1)=\mu(I_3)=\frac{\lambda}{4}] = \frac{2}{\lambda}(m_{k-1}(I_1)+m_{k-1}(I_3)) 
        \leq \frac{2\cdot 1}{\lambda}\cdot\frac{4\cdot 2}{\lambda}\max\{m_{k-2}(I_{1,1})+m_{k-2}(I_{1,3}),m_{k-2}(I_{3,1})+m_{k-2}(I_{3,3})\} = \cdots = \frac{2^1\cdot 1}{\lambda}\cdot
        \frac{2^2\cdot 2}{\lambda}\cdots\frac{2^k\cdot 2}{\lambda}\cdot 2\sup\limits_{-1<x<1}|f(x)| \leq \frac{2^{\frac{k(1+k)}{2}}}{\lambda^{k}}\cdot 2^k \leq \frac{2^{\frac{k(1+k)}{2}}}{\lambda^{k}}
        \cdot k^k$, $(k\leq 2)$ 
\end{enumerate}
\end{sol*}
\end{minipage}}
\newline\newline
\fbox{\begin{minipage}{0.95\textwidth}
\begin{figure}[H]
    \centering
    \includegraphics[width=0.32\textwidth]{"./Figures/250911_1.png"}
\end{figure}
\begin{enumerate}
    \item[c)] \textcolor{silkred}{(Wait to prove)}
\end{enumerate}
\end{minipage}}

\begin{prb}
    (C5S3Q7) \\
    Show that if a function $f$ is defined and differentiable on an open interval $I$ and $[a, b] \subset I$, then
    \begin{enumerate}
        \item[a)] the function $f'(x)$ (even if it is not continuous!) assumes on $[a, b]$ all the values between $f'(a)$ and $f'(b)$ (the theorem of Darboux);
        \item[b)] if $f''(x)$ also exists in $(a, b)$, then there is a point $\xi \in (a, b)$ such that $f'(b) - f'(a) = f''(\xi)(b-a)$.
    \end{enumerate}
\end{prb}
\fbox{\begin{minipage}{0.95\textwidth}
\begin{qsol*}
\begin{enumerate}
    \item[a)] (Fermat's Lemma + Weierstrass Maximum Theorem) /Constructive Proof/ \\
    Let $y$ be an arbitrary number between $f'(a)$ and $f'(b)$, and consider the function $F(x) = f(x) - yx$. Then $F'(x) = f'(x) - y$, and $F'(a) = f'(a) - y$, 
    $F'(b) = f'(b) - y$. $F^{\prime} (a)\cdot F^{\prime} (b)<0$.
    \item[b)] (Lagrange's MVT) \\
    \fbox{\begin{minipage}{0.9\textwidth}
        Lemma. $f^{\prime}$ doesn't have the 1st kind of discontinuity.
    \end{minipage}}
    The proof for Lemma is very easy, just choose a small neibourhood of discontinuity and using a). \\
    Apply Lagrange's MVT to $f'$ on $[a,b]$, we could get the result.
\end{enumerate}
\end{qsol*}
\end{minipage}}

\begin{prb}
    (C5S3Q8) \\
    A function $f(x)$ may be differentiable on the entire real line, without having a continuous derivative $f'(x)$
    \begin{enumerate}
        \item[a)] Show that $f'(x)$ can have only discontinuities of second kind.
        \item[b)] Find the flaw in the following "proof" that $f'(x)$ is continuous: \\
        Proof. Let $x_0$ be an arbitrary point on $\mathbb{R}$ and $f'(x_0)$ the derivative of $f$ at the point $x_0$. By definition of the derivative and Lagrange's theorem
        $f'(x_0) = \lim_{x \to x_0} \frac{f(x) - f(x_0)}{x - x_0} = \lim_{x \to x_0} f'(\xi) = \lim_{\xi \to x_0} f'(\xi)$,
        where $\xi$ is a point between $x_0$ and $x$ and therefore tends to $x_0$ as $x \to x_0$.
    \end{enumerate}
\end{prb}
\fbox{\begin{minipage}{0.95\textwidth}
\begin{qsol*}
\begin{enumerate}
    \item[a)] Proof is aformentioned.
    \item[b)] Lagrange's MVT requires the contimuity of function on closed interval.
\end{enumerate}
\end{qsol*}
\end{minipage}}

\begin{prb}
    (C5S3Q9) \\
    Let $f$ be twice differentiable on an interval $I$. Let $M_0 = \sup_{x \in I} |f(x)|$, $M_1 = \sup_{x \in I} |f'(x)|$ and $M_2 = \sup_{x \in I} |f''(x)|$. Show that
    \begin{enumerate}
        \item[a)] if $I = [-a, a]$, then $|f'(x)| \leq \frac{M_0}{a} + \frac{x^2 + a^2}{2a} M_2$;
        \item[b)] $\begin{cases} M_1 \leq 2\sqrt{M_0M_2}, & \text{if the length of } I \text{ is not less than } 2\sqrt{M_0/M_2}, \\ M_1 \leq \sqrt{2M_0M_2}, & \text{if } I = \mathbb{R}; \end{cases}$
        \item[c)] the numbers 2 and $\sqrt{2}$ in part b) cannot be replaced by smaller numbers;
        \item[d)] if $f$ is differentiable $p$ times on $\mathbb{R}$ and the quantities $M_0$ and $M_p = \sup_{x \in \mathbb{R}} |f^{(p)}(x)|$ are finite, then the quantities $M_k = \sup_{x \in \mathbb{R}} 
        |f^{(k)}(x)|$, $1 \leq k < p$, are also finite and $M_k \leq 2^{k(p-k)/2} M_0^{1-k/p} M_p^{k/p}$.
    \end{enumerate}
    Hint: Use Exercises 9b) and mathematical induction.
\end{prb}
\fbox{\begin{minipage}{0.95\textwidth}
\begin{sol*}
\begin{enumerate}
    \item[a)] Using Taylor's expansion: $f(x)=f(x_0)+f'(x_0)(x-x_0)+\frac{f''(\xi)}{2}(x-x_0)^2$ at $x=a$ and $x=-a$: \\
    $f(a)=f(x_0)+f'(x_0)(a-x_0)+\frac{f''(\xi_1)}{2}(a-x_0)^2$, $f(-a)=f(x_0)+f'(-a)(-a-x_0)+\frac{f''(\xi_2)}{2}(-a-x_0)^2$. \\
    Then $f'(x_0)=\frac{f(a)-f(-a)}{2a}-\frac{1}{4a}f''(\xi_1)(a-x_0)^2+\frac{1}{4a}f''(\xi_2)(a+x_0)^2$, using triangle inequality gives: $|f'(x)|\leq \frac{M_0}{a}+\frac{M_2}{4a}(a-x_0)^2+\frac{M_2}{4a}(a+x_0)^2
    =\frac{M_0}{a}+\frac{M_2}{2a}(a^2+x_0^2)$. 
    \item[b)] i) Consider $I=[a,b]$, the method is absolutely same as a). \\
    ii) Perform the analysis again for some intervals $[x - a, x)$, $(x, x + a]$ and we recover a similar result as above: $f(x - a) = f(x) + f'(x)(-a) + \frac{f''(c_1)}{2}(a)^2$, $f(x + a) = f(x) + f'(x)(a) + 
    \frac{f''(c_2)}{2}(a)^2$. \\
    Subtracting the two Taylor theorem estimates we obtain that:
    $f'(x) = \frac{f(x + a) - f(x - a)}{2a} - \frac{f''(c_1) + f''(c_2)}{4a}a^2$
    and we have that by triangle inequality:
    $|f'(x)| \leq \frac{M_0}{a} + \frac{M_2a}{2}$
    Since
    $|f'(x)| \leq \delta(a) \quad \forall a \in \mathbb{R}$
    it must be true that since $|f'(x)| \leq M_1$ for all $a$ as well, we must have that
    $M_1 \leq \inf_a \delta(a)$
    for if the opposite happens, then there will be some subset of the reals where the function violates the first inequality.
    We compute the infimum as follows:
    $\delta(a) = \frac{M_0}{a} + \frac{a}{2}M_2 \geq \frac{M_0}{a} + \frac{a}{2}M_2 \geq 2\sqrt{\frac{M_0M_2}{2}} = \sqrt{2M_0M_2}$
    and the proposition is proven. The infimum happens when $a = \sqrt{\frac{2M_0}{M_2}}$.
    \item[c)] Example: $f(x)=2x^2-1-a^2,\quad I=[a,1]$, $\forall a \in [0,1)$. 
    \item[d)] [Induction] \\
    $p=2$, $M_1\leq 2^{\frac{1}{2}}M_0^{\frac{1}{2}}M_2^{\frac{1}{2}}=\sqrt{2M_{0M_2}}$ \checkmark;  \\
    $p=2$, suppose $p$ holds $\forall 1\leq k\leq p-1$: \underline{$M_k\leq 2^{\frac{k(p-K)}{2}}M_0^{1-\frac{k}{p}}M_p^{\frac{k}{p}}$}. Then for $p+1$: 
    first, using MVT(Taylor's expansion is allright)    
\end{enumerate}
\end{sol*}
\end{minipage}}
\newline\newline
\fbox{\begin{minipage}{0.95\textwidth}
    $\begin{cases} &f^{(p-1)}(x+h)=f^{(p-1)}(x)+hf^{p}(x)+\frac{h^2}{2}f^{(p+1)}(\xi_1) \\ &f^{(p-1)}(x+h)=f^{(p-1)}(x)+hf^{p}(x)+\frac{h^2}{2}f^{(p+1)}(\xi_1) \end{cases}$, \\
    then $f^{p}(x) = \frac{1}{2h}(f^{(p-1)}(x+h)-f^{(p-1)}(x-h))-\frac{h}{4}[f^{(p+1)}(\xi_1)-f^{(p+1)}(\xi_2)]$, then $M_p\leq \frac{1}{h}M_{p-1}+\frac{h}{2}M_{p+1}=\left[
    h=\sqrt{2\frac{M_{p-1}}{M_{p+1}}} \right]=\sqrt{2M_{p-1}M_{p+1}}$, i.e. \underline{$M_p\leq 2^{\frac{1}{2}}M_{p-1}^{\frac{1}{2}}M_{p+1}^{\frac{1}{2}}$}. \\
    Using induction hypothesis, we have \underline{$M_{p-1}\leq 2^{\frac{p-1}{2}}M_0^{\frac{1}{p}}M_p^{\frac{1-p}{p}}$}, conbining 2 of them, then \underline{$M_p\leq 2^{\frac{p}{2}}M_0^
    {\frac{1}{p+1}}M_{p+1}^{\frac{p}{p+1}}$}, then we could verify the condition for $p+1$. 
\end{minipage}}

\begin{prb}
    (C5S3Q10) \\
    Show that if a function $f$ has derivatives up to order $n + 1$ inclusive at a point $x_0$ and $f^{(n+1)}(x_0) \neq 0$, then in the Lagrange form of the remainder in Taylor's formula
    $r_n(x_0; x) = \frac{1}{n!} f^{(n)} \left( x_0 + \theta (x - x_0) \right) (x - x_0)^n$,
    where $0 < \theta < 1$ and the quantity $\theta = \theta(x)$ tends to $\frac{1}{n+1}$ as $x \to x_0$.
\end{prb}
\fbox{\begin{minipage}{0.95\textwidth}
\begin{qsol*}
    Just compare $f(x)=\sum_{k=0}^{n-1}\frac{f^{(k)}(x_0)}{k!}(x-x_0)^k+\frac{f^{(n)}(x_0+\theta_x(x-x_0))}{n!}(x-x_0)^n$ and $f(x)=\sum_{k=0}^{n-1}\frac{f^{(k)}(x_0)}{k!}(x-x_0)^k+\frac{f^{(n)}(x_0)}{n!}(x-x_0)^n
    +\frac{f^{(n+1)}(x_0+\phi_x(x-x_0))}{(n+1)!}(x-x_0)^{n+1}$, then we get $\frac{f^{(n)}(x_0+\theta_x(x-x_0))}{n!} =\frac{f^{(n)}(x_0)}{n!} +\frac{f^{(n+1)}(x_0+\phi_x(x-x_0))}{(n+1)!}(x-x_0)$ i.e. $\frac{f^{(n)}
    (x_0+\theta_x(x-x_0))-f^{(n)}(x_0)}{x-x_0}=\frac{f^{(n+1)}(x_0+\phi_x(x-x_0))}{n+1}$, then $\theta_x=\frac{1}{n+1}\cdot\frac{f^{(n+1)}(x_0+\phi_x(x-x_0))}{ \dfrac{f^{(n)}(x_0+h_x)-f^{(n)}(x_0)}{h_x} }$, where 
    $h_x=\theta_x(x-x_0)$. Thus $\lim_{X \to x_0} \theta_x = \frac{1}{n+1}$. 
\end{qsol*}
\end{minipage}}

\begin{prb}
    (C5S4Q1) \\
    Let $x = (x_1, \ldots, x_n)$ and $\alpha = (\alpha_1, \ldots, \alpha_n)$, where $x_i \geq 0$, $\alpha_i > 0$ for $i = 1, \ldots, n$ and $\sum_{i=1}^n \alpha_i = 1$. For any number $t \neq 0$ we consider the mean of order $t$ of the numbers $x_1, \ldots, x_n$ with weights $\alpha_i$: $M_t(x, \alpha) = \left( \sum_{i=1}^n \alpha_i x_i^t \right)^{1/t}$.
    In particular, when $\alpha_1 = \cdots = \alpha_n = \frac{1}{n}$, we obtain the harmonic, arithmetic, and quadratic means for $t = -1, 1, 2$ respectively.
    Show that
    \begin{enumerate}
        \item[a)] $\lim_{t \to 0} M_t(x, \alpha) = x_1^{\alpha_1} \cdots x_n^{\alpha_n}$, that is, in the limit one can obtain the geometric mean;
        \item[b)] $\lim_{t \to +\infty} M_t(x, \alpha) = \max_{1 \leq i \leq n} x_i$;
        \item[c)] $\lim_{t \to -\infty} M_t(x, \alpha) = \min_{1 \leq i \leq n} x_i$;
        \item[d)] $M_t(x, \alpha)$ is a nondecreasing function of $t$ on $\mathbb{R}$ and is strictly increasing if $n > 1$ and the numbers $x_i$ are all nonzero.
    \end{enumerate}
\end{prb}
\fbox{\begin{minipage}{0.95\textwidth}
\begin{sol*}
\begin{enumerate}
    \item[a)] Attempt 1. (Convacity of $\ln x$) $\lim_{t \to 0} (\sum\limits_{i=1}^{n}\alpha_i x_i^t)^{\frac{1}{t}} = \lim_{t \to 0} e^{\frac{1}{t}\ln(\sum\limits_{i=1}^{n}\alpha_i x_i^t)} \geq \lim_{t \to 0} e^{\frac{1}{t}\sum\limits_{i=1}^{n} \alpha_i\ln x_i^n} = \prod_{i=1}^{n} x_i^{\alpha_i}$ \\
    Attempt 2. (L'Hôpital) $\lim_{t \to 0} (\sum\limits_{i=1}^{n}\alpha_i x_i^t)^{\frac{1}{t}} = \lim_{t \to 0} e^{\frac{1}{t}\ln(\sum\limits_{i=1}^{n}\alpha_i x_i^t)} = e^{\lim\limits_{t \to 0} \frac{\ln(\sum\limits_{i=1}^{n}\alpha_i x_i^t)}{t}} = e^{\lim\limits_{t \to 0} \frac{\sum\limits_{i=1}^{n}\alpha_i x_i^t \ln x_i}{\sum\limits_{i=1}^{n}
    \alpha_i x_i^t}} = e^{\sum\limits_{i=1}^{n}\alpha_i \ln x_i} = \prod_{i=1}^{n} x_i^{\alpha_i}$ 
    \item[b)] Let $\max_{1 \leq i \leq n}\{x_i\} = x_k$, then $\lim_{t \to \infty} (\alpha_k x_k^t)^{\frac{1}{t}} \leq \lim_{t \to \infty} M_t(x,\alpha) \leq \lim_{t \to \infty} ((\sum\limits_{i=1}^{n} \alpha_i)x_k^t)^{\frac{1}{t}}$. 
    \item[c)] Let $\min_{1 \leq i \leq n}\{x_i\} = x_l$, then $\lim_{t \to \infty} (\alpha_k x_l^t)^{\frac{1}{t}} \geq \lim_{t \to \infty} M_t(x,\alpha) \geq \lim_{t \to \infty} ((\sum\limits_{i=1}^{n} \alpha_i)x_l^t)^{\frac{1}{t}}$. 
    \item[d)] /Jensen's Inequality/ Notice that $y=x^m,\ x\geq 0$ is convex,then \\
    $(\sum_{i=1}^{n}\alpha_i x_i^s)^{\frac{t}{s}} \leq \sum_{i=1}^{n}\alpha_i (x_i^s)^{\frac{t}{s}} = \sum_{i=1}^{n}\alpha_i x_i^t \implies (\sum_{i=1}^{n} \alpha_i x_i^s)^{\frac{1}{s}} \leq (\sum_{i=1}^{n} \alpha_i x_i^t)^{\frac{1}{t}}$. Then $\forall s \leq t \in \mathbb{R}$, $M_s(x,\alpha) \leq M_t(x,\alpha)$. "$=$" i.f.f. $\forall x_i,\ i \in \{1,2,\cdots,n\}$, 
    $x_i$ all equals \textcolor{silkred}{(It's waiting to prove)}.
\end{enumerate}
\end{sol*}
\end{minipage}}

\begin{prb}
    (C5S4Q2) \\
    Show that $|1 + x|^p \geq 1 + px + c_p \varphi_p(x)$, where $c_p$ is a constant depending only on $p$,
    $\varphi_p(x) = \begin{cases} &|x|^2 \text{ for } |x| \leq 1, \\ &|x|^p \text{ for } |x| > 1, \end{cases}$, $\text{if } 1 < p \leq 2$, 
    and $\varphi_p(x) = |x|^p$ on $\mathbb{R}$ if $2 < p$.
\end{prb}
\fbox{\begin{minipage}{0.95\textwidth}
\begin{sol*}
\fbox{\begin{minipage}{0.75\textwidth}
    Lemma. $(1+x)^p \geq 1+px \quad \forall x\neq 0,\ p>1$.
\end{minipage}}
1$^\circ$. $1<p\leq 2$, $|x|\leq 1$. /Cauchy's MVT/ \\
\begin{equation*}
    \begin{split}
        \frac{(1+x)^{p-(1+px)}}{x^2} &= \frac{(p(1+x_1)^{p-1}-p)}{2x_1} \quad \exists x_i \in (-1,1)/\{0\} \\
        &= \frac{p(p-1)(1+x_2)^{p-2}}{2} \quad \exists x_2 \in (-x_1,x_1)/\{0\} \\
        &\geq \frac{p(p-1)}{2}
    \end{split}
\end{equation*}
then $(1+x)^p \geq 1+px+\frac{p(p-1)}{2}x^2$. \\
2$^\circ$. $1<p\leq 2$, $|x|>1$. \\
Let $g(x) = \frac{(1+x)^p - 1 - px}{|x|^p}$ is continuous and $g(x) \leq 0$ on $(-\infty,-1]\cup[1,\infty)$. $g(1)=\lim_{x \to 1^+} g(x)=\frac{2^p-(1+p)}{1}>0$, $g(-1)=\lim_{x \to 1^-} g(x)=\frac{p-1}{1}>0$, 
$\lim_{x \to +\infty} g(x)=\lim_{x \to -\infty} g(x)=1$. Then $g(x)>0$ on $(-\infty,-1)\cup(1,\infty)$, i.e. $\exists c_{p}>0$ s.t. $g(x)=\frac{(1+x)^p - 1 - px}{|x|^p}\geq c_p$, i.e. $(1+x)^p \geq 1 + px + c_p |x|^p$. \\
\end{sol*}
\end{minipage}}
\newline\newline
\fbox{\begin{minipage}{0.95\textwidth}
3$^\circ$. $p>2,\ x \in \mathbb{R}$ \\
Analysis is similiar with 2$^\circ$. Let $g(x) = \frac{(1+x)^p - 1 - px}{|x|^p}$ and analyze its boundary. (Maybe L'Hospital Rule is used in process.)
\end{minipage}}

\begin{prb}
    (C5S4Q3) \\
    Verify that $\cos x < (\frac{\sin x}{x})^3$, for $0<|x|<\frac{\pi}{2}$.
\end{prb}
\fbox{\begin{minipage}{0.95\textwidth}
\begin{qsol*}
Let $f(x)=\sin x (\cos x)^{-\frac{1}{3}}-x$, $x \in (0,\frac{\pi}{2})$ ($x \in (-\frac{\pi}{2},0)$ is same). 
\begin{equation*}
    \begin{split}
        f'(x) &= \cos^{\frac{2}{3}}x+\frac{1}{3}\sin^{2}x\cos^{-\frac{4}{3}}x-1 \\
        &= \cos^{\frac{2}{3}}x+\frac{1}{3}\cos^{-\frac{4}{3}}x-\frac{1}{3}\cos^{\frac{2}{3}}x-1 \\
        &= (\frac{1}{3}\cos^{-\frac{4}{3}}x-\frac{2}{3}\cos^{-\frac{2}{3}}x+\frac{1}{3})+\frac{2}{3}(\cos^{\frac{2}{3}}x-2+\cos^{-\frac{2}{3}}x) \\
        &= \frac{1}{3}(\cos^{-\frac{2}{3}}x-1)^2+\frac{1}{3}(\cos^{\frac{1}{3}}x-\cos^{-\frac{1}{3}}x)^2 \geq 0
    \end{split}
\end{equation*}
\end{qsol*}
\end{minipage}}

\begin{prb}
    (C5S4Q5) \\
    Show that if $f \in C\cap (a, b)$ and the inequality
    $f\left(\frac{x_1 + x_2}{2}\right) \leq \frac{f(x_1) + f(x_2)}{2}$
    holds for any points $x_1, x_2 \in (a, b)$, then the function $f$ is convex on $(a, b)$.
\end{prb}
\fbox{\begin{minipage}{0.95\textwidth}
\begin{sol*}
    Main method is [(Cauchy's Forward or Backward) Induction], first prove 2 lemma: \\
    \fbox{\begin{minipage}{0.9\textwidth}
        Lemma 1. $f(\frac{x_1 + x_2 +\cdots+ x_{2^n}}{2^n}) \leq \frac{f(x_1)+f(x_2)+\cdots+f(x_{2^n})}{2^n}$.
    \end{minipage}}
    \newline
    [Induction] \\
    \fbox{\begin{minipage}{0.9\textwidth}
        Lemma 2. If $f(\frac{x_1 +\cdots+ x_{n+1} }{n+1}) \leq \frac{f(x_1)+\cdots+f(x_{n+1})}{n+1}$ holds, then $f(\frac{x_1 +\cdots+ x_{n}}{n}) 
        \leq \frac{f(x_1)+\cdots+f(x_{n})}{n}$.
    \end{minipage}}
    \newline
    Then $\forall k \in \mathbb{Z}^+$, $f(\frac{x_1 +\cdots+ x_{k}}{k}) \leq \frac{f(x_1)+\cdots+f(x_{k})}{k},\ \forall x_i \in (a,b)$. \\
    Let $k=q$, $\forall 0<p<q \in \mathbb{Z}^+$, let $x_1=\cdots=x_p=x_p^{\prime},\ x_{p+1}=\cdots=x_k=x_2^{\prime}$, \\
    then $f(\frac{p}{q}x_1^{\prime}+(1-\frac{p}{q})x_2^{\prime}) \leq \frac{p}{q}f(x_1^{\prime})+(1-\frac{p}{q})f(x_2^{\prime})$. \\
    Since $f \in C(a,b)$, then $\forall \lambda \in (0,1)$, $f(\lambda x_1 + (1-\lambda)x_2) \leq \lambda f(x_1) + (1-\lambda)f(x_2)$, i.e. 
    $f$ is convex on $(a,b)$.
\end{sol*}
\end{minipage}}

\begin{prb}
    (C5S4Q6) \\
    Shaw that 
    \begin{enumerate}
        \item[a)] if a convex function $f : \mathbb{R} \to \mathbb{R}$ is bounded, it is constant;
        \item[b)] if $\lim\limits_{x \to -\infty} \frac{f(x)}{x} = \lim\limits_{x \to +\infty} \frac{f(x)}{x} = 0$, for a convex function $f : \mathbb{R} 
        \to \mathbb{R}$, then $f$ is constant.
        \item[c)] for any convex function $f$ defined on an open interval $a < x < +\infty$ (or $-\infty < x < a$), the ratio $\frac{f(x)}{x}$ tends to a 
        finite limit or to infinity as $x$ tends to infinity in the domain of definition of the function.
    \end{enumerate}
\end{prb}
\fbox{\begin{minipage}{0.95\textwidth}
\begin{sol*}
\begin{enumerate}
    \item[b)] 
    \fbox{\begin{minipage}{0.7\textwidth}
        Lemma. If $f:D\to\mathbb{R}$ is convex, then $\forall x_0 \leq x \leq y \in D$, satisfy $\frac{f(x)-f(x_0)}{x-x_0} \leq \frac{f(y)-f(x_0)}{y-x_0}$.
    \end{minipage}}
    Proof for lemma: let $f(x)=f(\alpha_1 x_0 + \alpha_2 y) \leq \alpha_1 f(x_0) + \alpha_2 f(y)$ and $x-x_0=\alpha_2(y-x_0)$. \\
    Using Lemma, let $x_0=0$, then $\frac{f(y)-f(0)}{y} \leq \frac{f(x)-f(0)}{x}$, $y \leq x,\ x,\ y\neq 0$ \\
    $\lim_{y \to -\infty} \frac{f(y)-f(0)}{y} = 0 \geq \frac{f(x)-f(0)}{x}$, $\forall  0<x \leq y$, then $f(x) \leq f(0)$, $x>0$, \\
    $\lim_{x \to -\infty} \frac{f(x)-f(0)}{x} = 0 \leq \frac{f(y)-f(0)}{y}$, $\forall  x<0 < y$, then $f(y) \leq f(0)$, $y>0$, \\
    Similiar with $x<0$, then $f(x)\equiv f(0)$, $\forall x \in \mathbb{R}$. 
    \item[a)] Method 1. By b), it's easy to prove. \\
    Method 2. [Contradiction] (Withoud b)) \\
    If not, $\exists x_1 < x_2 \in \mathbb{R}$ s.t. $f(x_1)<f(x_2)$ ($f(x_1)>f(x_2)$ is same), then $\forall \alpha_1 + \alpha_2 = 1, \exists x_3 > x_2$ s.t. 
    $x_2 = \alpha_1 x_1 + \alpha_2 x_3$, then $f(x_2) \leq \alpha_1 f(x_1) + \alpha_2 f(x_3)$, i.e. $\Pi$: $\begin{cases} &\frac{\alpha_1}{\alpha_2}[f(x_2)-f(x_1)] 
    \leq f(x_3)-f(x_2) \\ &\cdots\ \ \text{go on} \\ &\frac{\alpha_1}{\alpha_2}{f(x_n)-f(x_{n-1})} \leq f(x_{n+1})-f(x_n) \end{cases}$. Let $\frac{\alpha_1}{\alpha_2}>1$, 
    then $\left(\frac{\alpha_1}{\alpha_2}\right)^n[f(x_2)-f(x_1)] \leq f(x_{n+1})-f(x_n)$, i.e. $\exists \{x_n\}_{n=1}^{+\infty}\nearrow$, $f(x_{n+1}) \to +\infty$ as 
    $n \to +\infty$, contradicting the boundedness of $f$.
    \item[c)] Still using b). lemma, construct $g(x) \triangleq \frac{f(x)-f(a)}{x-a},\ x \in (a,+\infty)$ (define $f(a)=\lim_{x \to a^+} f(x)$). Then $g(x) \nearrow$, 
    $x>a$, with Weierstrass Theorem, $\lim_{x \to +\infty} g(x) = c \ (\leq \infty)$. For $f(x)$, $x \leq -a$, $g(x)$ is similar.
\end{enumerate}
\end{sol*}
\end{minipage}}

\begin{prb}
    (C5S4Q7) \\
    Show that if $f : (a, b)\rightarrow \mathbb{R}$ is a convex function, then
    \begin{enumerate}
        \item[a)] at any point $x \in (a, b)$ it has a left-hand derivative $f'_-$ and a right-hand derivative $f'_+$, defined as $f'_-(x) = 
        \lim\limits_{h \to -0} \frac{f(x + h) - f(x)}{h}$, $f'_+(x) = \lim\limits_{h \to +0} \frac{f(x + h) - f(x)}{h}$, and $f'_-(x) \leq f'_+(x)$;
        \item[b)] the inequality $f'_+(x_1) \leq f'_-(x_2)$ holds for $x_1, x_2 \in (a, b)$ and $x_1 < x_2$;
        \item[c)] the set of cusps of the graph of $f(x)$ (for which $f'-(x) \neq f'+(x)$) is at most countable.
    \end{enumerate}
\end{prb}
\fbox{\begin{minipage}{0.95\textwidth}
\begin{sol*}
\begin{enumerate}
    \item[a)] Method 1. [Proof by contradiction] Using definition, it's easy to get contradiction. \\
    Mathod 2. Using lemma in (C5S4Q6)b), direct result obtained.
    \item[b)] 
    \fbox{\begin{minipage}{0.9\textwidth}
        Lemma. (Developement in (C5S4Q6)b)) If $f:D\to\mathbb{R}$ is convex, then $\forall x \in D$, $\forall x, y: x \leq x_0 \leq y \in D$, s.t. 
        $\frac{f(x)-f(x_0)}{x-x_0} (\leq \frac{f(y)-f(x_0)}{y-x_0}) \leq \frac{f(y)-f(x)}{y-x}$.
    \end{minipage}}
\end{enumerate}     
\end{sol*}
\end{minipage}}
\newline\newline
\fbox{\begin{minipage}{0.95\textwidth}
\begin{enumerate}
    \item[b)] Proof for Lemma: similiar with (C5S4Q6)b). (using twice). \\
    $\forall x_1<x_2\in (a,b)$, $x_1<x_2$, $\exists h_+>0, h_-<0$, s.t. $x_1<x_1+h_+<x_2+h_-<x_2$. Then using twice of Lemma, we could have $\frac{f(x_1+h_+)-f(x_1)}{h_+} 
    \leq \frac{f(x_2+h_-)-f(x_1+h_+)}{x_2+h_--(x_1+h_+)} \leq \frac{f(x_2)-f(x_2+h_-)}{-h_-} \implies f_+^{\prime} (x_1) \leq f_-^{\prime} (x_2)$.
    \item[c)] The proof is similiar with the one of "The first kind of discontinuities of a monotone function are at most countable". \\
    Proof idea and recall: $\forall x_0 \in \{x:f_{-}^{\prime} (x) \neq f_{+}^{\prime} (x)\} = \{x:f_{-}^{\prime} (x) < f_{+}^{\prime} (x)\}$, $\exists$ an open interval 
    $(f_{-}^{\prime} (x_0), f_{+}^{\prime} (x_0))$ corresponding to $x_0$. And due to "The number of open intervals on the real number axis that do not intersect each other 
    is at most countable", then problem is solved. 
\end{enumerate}
\end{minipage}}

\begin{prb}
    (C5S4Q8) \\
    The Legendre transform of a function $f : I \to \mathbb{R}$ defined on an interval $I \subseteq \mathbb{R}$ is the function $f^*(t) = \sup\limits_{x \in I} (t x - f(x))$. \\
    Show that
    \begin{enumerate}
        \item[a)] The set $I^*$ of values of $t \in \mathbb{R}$ for which $f^*(t) \in \mathbb{R}$ (that is, $f^*(t) \neq \infty$) is either empty or consists of a single point, 
        or is an interval of the line, and in this last case the function $f^*(t)$ is convex on $I^*$.
        \item[b)] If $f$ is a convex function, then $I^* \neq \emptyset$, and for $f^* \in C(I^*)$ $(f^*)^* = \sup\limits_{t \in I^*} \left( x t - f^*(t) \right) = f(x)$ for any $x \in I$. 
        Thus the Legendre transform of a convex function is involutive, (its square is the identity transform).
        \item[c)] The following inequality holds: $x t \leq f(x) + f^*(t)$ for $x \in I$ and $t \in I^*$.
        \item[d)] When $f$ is a convex differentiable function, $f^*(t) = t x_t - f(x_t)$, where $x_t$ is determined from the equation $t = f'(x)$. Use this relation to obtain 
        a geometric interpretation of the Legendre transform $f^*$ and its argument $t$, showing that the Legendre transform is a function defined on the set of tangents to the 
        graph of $f$.
        \item[e)] The Legendre transform of the function $f(x) = \frac{1}{\alpha} x^{\alpha}$ for $\alpha > 1$ and $x \geq 0$ is the function $f^*(t) = \frac{1}{\beta} t^{\beta}$, 
        where $t \geq 0$ and $\frac{1}{\alpha} + \frac{1}{\beta} = 1$. Taking account of c), use this fact to obtain Young's inequality, which we already know: $x t \leq 
        \frac{1}{\alpha} x^{\alpha} + \frac{1}{\beta} t^{\beta}$.
        \item[f)] The Legendre transform of the function $f(x) = e^{x}$ is the function $f^*(t) = t \ln \frac{t}{e}, t > 0$, and the inequality $x t \leq e^{x} + t \ln \frac{t}{e}$ 
        holds for $x \in \mathbb{R}$ and $t > 0$.
    \end{enumerate}
\end{prb}
\fbox{\begin{minipage}{0.95\textwidth}
\begin{sol*}
\begin{enumerate}
    \item[a)] Condition 1. $I$ is bounded, but $f$ is unbounded, or $I$ and $f$ are unbounded, but $|f|$ grows faster then $px$. \\
    Condition 2. $I$ and $f$ are unbounded, and $f$ grows as fast as $px$, for some $p$. \\
    Condition 3. $I$ is 
\end{enumerate}    
\end{sol*}
\end{minipage}}

%\end{document}